% =====================================================================
% NOTE TO AUTHOR (read before compiling):
%   September 14, 2026 revision draft based on the rejected BlackboxNLP 2026
%   submission; no future venue has been selected. Drop THIS file
%   (acl_latex.tex) and custom.bib into the official ACL Overleaf template
%   (ships acl.sty + acl_natbib.bst), then compile with pdfLaTeX + BibTeX.
%   The draft remains in [review] mode with an anonymous author block.
%
%   DATA: stagewise, axis-targeted-control, and corpus numbers come from
%   out/results/summary_v3.json (plus-one-corrected run, 5 canonical seeds,
%   4 extension seeds, 140 prompts) and matched_control_stagewise.json;
%   audit numbers come from ai_consensus_validation_v3.json. Derived stats via experiments/monte_carlo_correction.py,
%   derive_paper_stats.py, between_condition_stats.py, stage_condition_gaps.py,
%   multiple_testing.py, threshold_robustness.py, audit_cluster_bootstrap.py,
%   dpo_pairwise.py, dpo_phrase_robustness.py, and raw_output_robustness.py.
%   Primary inference = prompt-clustered permutation tests (in the artifact).
%   COMPLETED DIAGNOSTIC: two independent blinded AI passes over 120 sentences,
%   with three disagreements adjudicated by Claude Opus 4.8. This is explicitly
%   AI-only and is not human validation. PENDING: causal leave-cluster-out,
%   Stage-2 Dolmino index, and independent human validation. No prospective
%   experiment or human-label result is represented as complete in this draft.
% =====================================================================

\documentclass[11pt]{article}

\usepackage[review]{acl}

\usepackage{times}
\usepackage{latexsym}
\usepackage[T1]{fontenc}
\usepackage[utf8]{inputenc}
\usepackage{microtype}
\usepackage{booktabs}
\usepackage{graphicx}
\usepackage{amsmath}

% Absorb small overfulls from long \texttt{} identifiers.
\emergencystretch=3em

\title{Measuring Advice-Style Language in OLMo-2:\\
Effects of Detection Rules and Response Length}

% === Anonymous author block for review (current) ===
\author{Anonymous Submission}

\begin{document}
\maketitle

\begin{abstract}
Advice responses often contain phrases of acknowledgment, validation,
professional guidance, and list organization. Their measured prevalence depends
on what counts as a match and how much text a model produces. We ask: how do
detection rules and response length change
conclusions about advice-style language across OLMo-2 post-training? We apply a
fixed lexicon--embedding detector to a 140-prompt battery spanning five
conditions and public 1B and 7B checkpoints, then audit its components and
output windows. A 120-sentence AI-label diagnostic finds low union-detector
precision, with false positives from the embedding component;
these labels are not human validation. In the recovered 700-response final-7B
cache, a 64-token window reverses the advice versus emotional-control contrast
for the category named
``empathy-opener.'' Checkpoint comparisons also depend on interface and detector
scope: Base-to-SFT changes confound training with prompt serialization, while
same-interface SFT-to-DPO trajectories differ across marker categories. Finally,
chosen-side referral/disclaimer and structuring phrase enrichment in roughly
378,000 DPO pairs per model shrinks sharply among responses within 20 words in
length. The study therefore measures detector-defined language patterns and
their sensitivity; it does not estimate semantic empathy, referral quality, or
causal effects of training examples.
\end{abstract}

\section{Introduction}

Consider two answers to ``I've been feeling anxious and cannot sleep---what
should I do?'' One begins, ``That sounds really difficult. Here are a few
steps,'' while another gives the same information without either phrase. A
substring rule, a sentence-embedding rule, and a fixed 64-token window may
classify these answers differently. If one condition also elicits longer
answers, the detector has more sentences in which to find a match. The resulting
rate is therefore a property of the responses \emph{and} the measurement rule.

This paper asks one question: \emph{How do detection rules and response length
change conclusions about advice-style language across OLMo-2 post-training?}
A \textbf{marker} is a specified textual feature, such as one of the phrases in
our lexicon. A \textbf{detector} is the fixed algorithm that predicts whether a
marker category occurs. A detector output is a \textbf{proxy} only when used as
an indirect stand-in for a broader construct such as empathy. We use these
terms consistently below. In particular, the category named
\emph{empathy-opener} in the original artifact is searched anywhere in an
answer, and \emph{referral/disclaimer} combines referrals, disclaimers, and
professional-language mentions. Neither output is a semantic judgment.

OLMo-2 is useful for this measurement study because its final Instruct model is
the product of four public data-bearing stages:
pretraining, supervised fine-tuning (SFT), direct preference optimization
(DPO), and reinforcement learning with verifiable rewards (RLVR)
\citep{olmo2025olmo2}. OLMo-2 offers an unusually complete public chain of
post-training data and intermediate checkpoints. We evaluate the same fixed
detector across checkpoints and use exact phrase searches in available training
data as supporting context. These observations localize detector changes and
co-occurrences; they do not attribute an output to a training example.

\paragraph{Contributions.}
\begin{enumerate}\itemsep2pt
\item We show that conclusions about condition differences change across exact,
embedding, union, and fixed-window rules; the available AI audit identifies
substantial embedding-driven false positives but does not replace human labels.
\item Under the historical detector, same-interface SFT-to-DPO changes differ
by category: validation gaps increase at both model sizes, while structuring
gaps contract; Base-to-SFT changes remain confounded with prompt serialization.
\item We show that aggregate DPO phrase differences are strongly associated
with within-pair response length: referral/disclaimer differences approach zero
and structuring differences shrink among similarly sized responses.
\end{enumerate}

\section{Related Work}

\paragraph{Behavioral tests and checkpoint trajectories.}
Our axis-targeted prompt battery follows the behavioral-testing principle that
targeted capabilities require controlled tests beyond aggregate accuracy
\citep{ribeiro2020checklist}. Pythia makes dense public checkpoints available
for training-dynamics research \citep{biderman2023pythia}; checkpoint studies
have tracked how loss and downstream behavior change during pretraining
\citep{xia2023trajectories}. We instead compare named post-training stages and
test changes in condition gaps rather than raw behavior rates.

Very recent work traces situation modeling and putative mentalizing across
OLMo-2 and Pythia stages while stress-testing construct validity
\citep{riviere2026developmental}. It is closest in model family and
developmental framing; our target is different (free-form advice markers and
their training-data co-occurrence), but its caution about mapping task outputs
to cognitive constructs directly motivates our proxy terminology.

\paragraph{Empathy and emotional-support measurement.}
Emotional-support research typically grounds support strategies in theory and
human annotation \citep{liu2021emotional}. Empathy research also warns that
different operational definitions can make results difficult to compare
\citep{lahnala2022critical}. Recent studies compare human and LLM judgments
directly \citep{welivita2026empathetic}; in mental-health settings, LLM judges
show reduced reliability for affective properties such as empathy
\citep{badawi2026trust}. These results make response-level human validation the
appropriate standard. Our AI-only audit diagnoses the rule; it does not meet
that standard.

\paragraph{Length in preference evaluation.}
Preference and reward-model data can encode length correlations
\citep{singhal2023long}, and length-controlled evaluation can change model
rankings \citep{dubois2024length}. Our contribution is narrower: we test how
response length changes marker-specific exact-match differences in the OLMo-2
DPO mixtures. The post-hoc length strata describe an association and are not a
causal adjustment.

\paragraph{Pretraining attribution and n-gram search.}
infini-gram \citep{liu2024infinigram} makes exact-substring counting over
web-scale corpora tractable, and OLMoTrace \citep{liu2025olmotrace} traces
model outputs back to pretraining text. Both are pretraining-only; we add
role-scoped search over the post-training layers, where our stagewise results
describe transition-level changes.

\paragraph{Open instruction and preference data.}
Tulu-3 \citep{lambert2024tulu3} is the recipe behind OLMo-2's post-training;
the SFT mixtures, DPO preference mixes, and RLVR sets are public, as is oasst1
\citep{kopf2023oasst1} (itself an internal subset of the Tulu-3 mixtures).
Public post-training data and checkpoints make the trace auditable.

\paragraph{The superficial alignment hypothesis.}
LIMA \citep{zhou2023lima} argues alignment largely surfaces a latent style
from few demonstrations; \citet{ghosh2024closer} interrogate this. Our
stagewise gaps give a marker-level version: each proxy is nonzero at base,
and post-trained checkpoints show larger advice-minus-factual detector gaps.

\paragraph{Causal data attribution.}
LESS uses optimizer-aware gradient similarity to select influential instruction
data \citep{xia2024less}; \citet{xiao2026probe} perform probe-based causal
attribution on OLMo-2 DPO and mitigate a behavior by editing top-attributed
preference pairs. We differ in kind: our trace is observational, spans multiple
public post-training checkpoints and datasets, and nominates clusters and stages
for such interventions; it does not establish causation.

\section{Measurement and Study Design}

\subsection{Operational categories and fixed detector}
\label{sec:method-detect}
A hand-seeded, inspectable lexicon specifies four marker categories using the
artifact's historical names: \emph{empathy-opener}, \emph{validation-phrase},
\emph{referral/disclaimer} (artifact key \texttt{disclaimer}), and
\emph{structuring-marker} (\texttt{structure}). The names label rules rather
than validated constructs. The empathy-opener rule searches every sentence, not
only an opening span. Referral/disclaimer is a broad combined rule that also
includes professional-language fragments such as ``consult a'' and ``medical
professional''; it is not a detector of referrals alone. A fifth historical
category, \texttt{crisis\_referral}, is excluded from the four-category analysis
after producing 0 true positives and 11 false positives in the audit.

The comparisons treat the archived detector outputs and primary threshold as
fixed rather than selecting rules after inspecting a contrast. The detector
takes the union of case-insensitive substring matching and an e5-small embedding rule
\citep{wang2024e5}: a response is positive when either an exact phrase appears
or any sentence reaches cosine similarity $\tau=0.82$ to a category exemplar.
We also examine $\tau=0.80$ and $0.84$. Longer responses supply more sentences
and therefore more opportunities for the embedding rule to fire. The complete
specification used to reproduce the July final-7B response analysis is in
Appendix~\ref{app:lexicon}; the raw-cache diagnostics preserve that version
rather than retrospectively improving it. Archived stage summaries are analyzed
as historical outputs. Because most stagewise raw responses are unavailable, we
do not claim to have re-executed every stage under a byte-identical detector.

\paragraph{Construction and pilot overlap.} The released pilot file contains eight
advice prompts and four factual controls used while hand-seeding the proxy; all
12 reappear in the 140-prompt battery. We therefore do not call the full battery
held out. A final-7B sensitivity excluding those 12 leaves the advice--factual
gaps similar (Appendix~\ref{app:response-robustness}). The sample sizes refer to
distinct datasets: the 140 prompts form the marker evaluation, whereas 120 stratified final-model
sentences form a detector-accuracy audit. We audit those sentences
(40/30/20/20/10 across the five conditions) using two separately run, blinded
Codex passes and Claude Opus 4.8 adjudication of their three disagreements.
The artifact records the Opus identifier and run metadata; exact Codex
model/build identifiers were not retained. This is an AI diagnostic, not human
validation. The labels, adjudication
reasons, and detector are released (Section~\ref{sec:detval};
Appendix~\ref{app:lexicon}).

\subsection{Prompt conditions and generation}
\label{sec:method-battery}
Five conditions target different aspects of advice-style responses:
\emph{advice} (60 personal or emotional help requests), \emph{factual} (30
knowledge questions), \emph{emotional non-advice} (20 descriptive prompts
about emotional topics without a help request), \emph{neutral advice} (20 how-to
requests without emotional stakes), and \emph{domain factual} (10 impersonal
clinical-information questions). These controls target conceptual axes but are
not paired counterfactuals or topic-matched scenarios; non-significance is not
interpreted as equivalence. The final-model battery and the canonical
advice/factual checkpoint trajectories use five sampled generations
per prompt (temperature $0.7$, top-$p$ $0.9$, 256 new tokens, distinct seeds).
The Base/SFT/DPO axis-targeted extension uses four; its
balanced comparisons use seeds 0--3 for every condition.

\subsection{Inference at the prompt level}
\label{sec:method-stats}
The clustering unit is the prompt: each prompt contributes its mean
detector-positive indicator over its four or five generations, and condition rates are means over
prompts. All between-condition inference uses prompt-clustered permutation
tests (10{,}000 permutations of prompt-level means; two-sided;
Monte Carlo $p=(b+1)/(B+1)$). Benjamini--Hochberg correction is applied
separately to the 24 final-condition contrasts, 28 advice/factual adjacent-stage
interactions, 24 axis-targeted stage interactions, and 8 primary DPO pairwise tests.
We do \emph{not} treat generations as independent Bernoulli trials:
all condition and stage tests operate on prompt-level generation means.

\subsection{Checkpoint evaluation and supporting corpus search}
\label{sec:method-stagewise}
The same evaluation runs on every public checkpoint: 1B
Base$\to$SFT$\to$DPO$\to$RLVR1$\to$Instruct; 7B
Base$\to$SFT$\to$DPO$\to$Instruct. Base checkpoints have no chat template and
use a plain Question/Answer format, whereas post-trained checkpoints use their
official chat template. Thus Base$\to$SFT is a stage-plus-interface
transition and cannot isolate SFT training. No serialization crossover was run;
we demote that transition throughout. The
stagewise grid runs the 90 advice$+$factual prompts at every checkpoint. All five
conditions run on the final models and, in the four-seed specificity extension,
at Base, SFT, and DPO.

Corpus searches play a supporting role. For each stage's available indexed data
we search the exact Stage-1 corpus
(\texttt{v4\_olmo-mix-1124\_llama}; novelty $=1-$verbatim coverage), each
model's exact SFT mixture (\texttt{tulu-3-sft-olmo-2-mixture-0225} for 1B;
\texttt{tulu-3-sft-olmo-2-mixture} for 7B) in DuckDB with assistant-only scoping
(866{,}138 conversations for 1B and 939{,}344 for 7B; oasst1 is not searched
separately since it is an internal subset), the DPO preference mixes
(pair-level chosen-only/rejected-only counts, exact McNemar tests, and phrase
occurrences per million response words), and the
RLVR sets. SFT recoverability compares seeded marker phrases against
length-matched topical controls from the same answers, per-million documents,
with a distinct-phrase leave-one-out. Stage-2 mid-training uses
Dolmino-Mix-1124, for which no public index was available; the corpus trace is
therefore incomplete.

\paragraph{Reproducibility.} Stagewise, condition, and corpus numbers derive
from \texttt{summary\_v3.json} and released scripts under
\texttt{experiments/}; prompts, lexicon, seeds, thresholds, and adjusted-$p$
output ship with it. Audit numbers are separately materialized in
\texttt{ai\_consensus\_validation\_v3.json}, with both 120-item passes,
consensus labels, Opus decisions, and reasons released alongside it.

\section{Results}
\label{sec:results}

We first audit the measurement rule, then report condition and checkpoint
comparisons under that rule, and finally examine the association between DPO
phrase differences and response length. This order matters because the audit
limits what the subsequent detector-positive rates can mean.

\subsection{Measurement audit: precision, components, and output windows}
\label{sec:detval}

A condition-stratified random audit samples 120 final-model sentences across
all five conditions (40 advice, 30 factual, 20 emotional non-advice, 20 neutral
advice, and 10 domain factual). The sample contains 56 sentences from the 1B
model and 64 from the 7B model, drawn from 97 distinct seed-0 responses. It was
not stratified on detector predictions, and multiple sentences can come from
one response. Two independent blinded Codex passes agree exactly on 117 of 120
sentences (97.5\%). Category agreement is 99.2--100\%; Cohen's $\kappa$ is
.885 for empathy expression, .663 for validation, 1.000 for referral, and .978
for structuring. Claude Opus 4.8 adjudicates the three disagreements. These are
AI-consensus labels under specific prompts and models, not human ground truth.

Against that consensus, the union detector has low precision and
category-dependent recall (Table~\ref{tab:detval}). The audit's referral label
is narrower than the detector: it requires an actual recommendation to seek a
qualified professional, while the historical referral/disclaimer rule also
matches disclaimers and professional-language mentions. Its 0.06 precision
therefore reflects both broad detector behavior and target misalignment; it
should not be read as the accuracy of a pure referral detector. Empathy and
validation recall are each 0.50. Errors vary by condition, so subtracting
condition rates does not guarantee cancellation. The detector overpredicts
positives in this sample, although the rates are not formal upper bounds because
recall is below one.

Exact substring matching has no false positives in these 120 sentences but
misses many AI-labeled positives. Its true-positive/false-positive/false-negative
counts are 2/0/2 for empathy expression, 1/0/1 for validation, 1/0/0 for
referral, and 9/0/21 for structuring. Every union false positive is introduced
by the embedding component. The union raises structuring recall from 0.30 to
0.87, with lower precision. Component choice is therefore part of the
measurement question.

The recovered 700-response cache permits a separate response-level diagnostic
for the final 7B checkpoint only. Advice responses average 244 tokens and
factual responses 187; the three targeted-control means are 251--255 tokens.
Restricting each response to its first 64 tokens changes the union detector's
advice-minus-emotional contrast for the category historically called
empathy-opener from $-.147$ to $+.253$. Validation remains positive, moving
from $.250$ to $.177$; referral/disclaimer moves from $.090$ to $.113$; and
structuring remains near zero, moving from $-.030$ to $-.033$. On full
responses, exact-only advice-minus-targeted-control contrasts are $.323$ for
empathy, $.153$ for validation, $.157$ for referral/disclaimer, and $.390$ for
structuring. These component and window checks qualify the final-7B result; raw
responses are unavailable for the other checkpoints, so they do not validate
the stagewise trajectories.

\begin{table}[t]
\centering
\small
\begin{tabular}{lcccc}
\toprule
category & P & R & F1 & tp/fp/fn \\
\midrule
empathy expr. & 0.22 & 0.50 & 0.31 & 2/7/2 \\
validation & 0.33 & 0.50 & 0.40 & 1/2/1 \\
ref./discl. & 0.06 & 1.00 & 0.12 & 1/15/0 \\
structuring & 0.49 & 0.87 & 0.63 & 26/27/4 \\
\bottomrule
\end{tabular}
\caption{Union-detector precision, recall, and F1 against the adjudicated
120-sentence AI consensus ($\tau=0.82$). This is a diagnostic against an AI
annotation construct, not human validation. The referral annotation is narrower
than the combined historical referral/disclaimer rule. Response-cluster
bootstrap intervals appear in Appendix~\ref{app:audit-bootstrap}.}
\label{tab:detval}
\end{table}

\subsection{Condition comparisons depend on the measurement rule}
\label{sec:decomposition}

Table~\ref{tab:conditions} gives mean prompt-level proxy-positive rates per condition on
the final models; Table~\ref{tab:contrasts} tests advice against each axis-targeted
control. Against plain factual controls every category separates sharply in both
models (all prompt-clustered $p\leq.0002$). For referral/disclaimer, the 1B
model has rates .900 on advice and .080 on factual prompts; the 7B model has
rates .913 and .060.

Targeted controls change the interpretation. At $\tau=.82$, the category
historically called empathy-opener fires more often on emotional non-advice than
advice: .790 versus .590 for 1B ($p=.010$), and .760 versus .617 for 7B
($p=.098$). As the preceding audit shows, the final-7B direction reverses in a
64-token window. Structuring is high on both advice forms: the personal-minus-
neutral advice difference is .000 for 1B ($p=1.0$) and $-.030$ for 7B
($p=.574$). These are tests for a difference, not equivalence tests.

The broad referral/disclaimer category is also common outside personal advice.
It fires on .740 of neutral-advice prompts for 1B and .720 for 7B, and on .820
of domain-factual prompts for both models. Its advice-minus-domain-factual
difference is small and not significant: .080 for 1B ($p=.352$) and .093 for
7B ($p=.159$). The validation-phrase union output is higher on advice than every
targeted control for both models (all unadjusted $p\leq.0044$; all BH-adjusted
$q\leq.0082$). These marginal profiles differ, but weak correlations do not
establish the absence of a latent response template. The comparisons are
detector-dependent, and the small, mixed 1B--7B differences do not support a
scaling claim.

\begin{table*}[t]
\centering
\small
\setlength{\tabcolsep}{5pt}
\begin{tabular}{llrrrrr}
\toprule
Marker proxy & Model & Advice & Factual & Emotional & Neutral adv. & Domain fact. \\
\midrule
empathy opener  & 1B & 0.590 & 0.040 & \textbf{0.790} & 0.020 & 0.340 \\
                & 7B & 0.617 & 0.027 & \textbf{0.760} & 0.070 & 0.360 \\
validation phr. & 1B & 0.363 & 0.007 & 0.120 & 0.000 & 0.040 \\
                & 7B & 0.383 & 0.013 & 0.130 & 0.030 & 0.060 \\
ref./discl.     & 1B & 0.900 & 0.080 & 0.510 & 0.740 & 0.820 \\
                & 7B & 0.913 & 0.060 & 0.540 & 0.720 & 0.820 \\
structuring     & 1B & 1.000 & 0.580 & 0.820 & \textbf{1.000} & 0.920 \\
                & 7B & 0.970 & 0.547 & 0.830 & \textbf{1.000} & 0.840 \\
\bottomrule
\end{tabular}
\caption{Mean prompt-level union-proxy rate per condition (final models, $\tau=0.82$;
each prompt averaged over 5 generations). Bold marks the cells that break the
``advice-specific'' reading: the empathy-opener proxy peaks on emotional
\emph{non-advice} prompts, and structuring is at ceiling on \emph{neutral} advice. The
audit sample shows substantial overprediction, suggesting upward-biased absolute
rates (Section~\ref{sec:detval}); \texttt{crisis\_referral} is
excluded (0 TP / 11 FP in the 120-sentence audit).}
\label{tab:conditions}
\end{table*}

\begin{table}[t]
\centering
\small
\setlength{\tabcolsep}{2pt}
\resizebox{\columnwidth}{!}{%
\begin{tabular}{llrrr}
\toprule
Marker & Model & $-$Emot. & $-$Neut.adv. & $-$Dom.fact. \\
\midrule
empathy    & 1B & $-.200$ (.010) & $.570$ (.0002) & $.250$ (.016) \\
           & 7B & $-.143$ (.098) & $.547$ (.0002) & $.257$ (.031) \\
validation & 1B & $.243$ (.003)  & $.363$ (.0002) & $.323$ (.003) \\
           & 7B & $.253$ (.003)  & $.353$ (.0003) & $.323$ (.004) \\
referral   & 1B & $.390$ (.0002) & $.160$ (.006) & $.080$ (.352) \\
           & 7B & $.373$ (.0002) & $.193$ (.001) & $.093$ (.159) \\
structuring & 1B & $.180$ (.0002) & $.000$ (1.0) & $.080$ (.018) \\
           & 7B & $.140$ (.004) & $-.030$ (.574) & $.130$ (.026) \\
\bottomrule
\end{tabular}%
}
\caption{Advice minus each axis-targeted control (prompt-clustered permutation $p$,
10k draws). Benjamini--Hochberg correction across the 24 displayed tests leaves
the validation-phrase result unchanged (all $q\leq.0082$). Non-significant
cells mean ``no detected difference,'' not equivalence.}
\label{tab:contrasts}
\end{table}

\paragraph{Marker co-occurrence.}
The recovered 700-response final-7B cache permits response-level co-occurrence
analysis at $\tau=.82$. Among advice responses, 4.7\% contain one proxy, 29.0\%
two, 40.7\% three, and 25.7\% all four (none contain zero). Referral/disclaimer
is present in 93.0\% of validation-positive advice responses, while pairwise phi
correlations range from $-.14$ to $.21$. The categories often co-occur, but
these weak correlations alone neither establish nor rule out a latent response
template. Condition-level joint distributions and all pairwise correlations are released
(Appendix~\ref{app:response-robustness}).

\subsection{Checkpoint contrasts under the historical detector}
\label{sec:stagewise}

Raw advice-side rates rising across stages would be ambiguous because post-training
could simply make \emph{all} outputs more cautious and listy. A more informative
operational quantity is the advice-minus-factual gap $\Delta_s$ at each stage
(Table~\ref{tab:stagewise}) and its \emph{change} across adjacent stages, a
stage$\times$condition difference-in-differences tested with a prompt-clustered
permutation test on per-prompt stage trajectories (Table~\ref{tab:interaction};
10k permutations, with 60 advice prompts and 30 factual prompts).

All eight Base-to-SFT/interface changes are positive and six survive BH
correction. For referral/disclaimer, the change is +.253 for the 1B model and
+.243 for the 7B model; for structuring it is +.337 and +.227, respectively.
Validation does not survive correction ($q=.137$ for 1B and $q=.121$ for 7B).
Every category is already detector-positive at Base. Because Base uses plain QA
and SFT uses the chat template, these changes cannot isolate training from
serialization.

The same-interface SFT-to-DPO transition differs by category. The validation
gap increases by +.213 for 1B ($q=.002$) and +.147 for 7B ($q=.021$). The 1B
referral/disclaimer gap increases by +.190 ($q=.002$), while its 7B change of
+.083 is not significant after correction ($q=.203$). The empathy-opener
changes are also not significant after correction. The structuring gap
contracts by $-.367$ for 1B and $-.277$ for 7B (both $q=.002$), because DPO
raises detector-positive structuring on factual answers while advice-side rates
are near ceiling. Aggregate preference-data co-occurrence cannot independently
explain these trajectories because its referral difference approaches zero
among similar-length pairs (Table~\ref{tab:dpo}).

No RLVR or final-merge gap change is significant (all raw $p\geq.168$). These
tests describe the outputs of the fixed detector. Because the audit finds
condition-dependent false positives and does not span checkpoints, they do not
by themselves establish changes in human-judged behavior prevalence.

\begin{table}[t]
\centering
\small
\setlength{\tabcolsep}{3pt}
\resizebox{\columnwidth}{!}{%
\begin{tabular}{llrr}
\toprule
Model & Marker proxy & Base$\to$SFT/interface & SFT$\to$DPO \\
\midrule
1B & empathy opener & $+0.190$ ($q=.004$) & $+0.117$ ($q=.121$) \\
1B & validation phrase & $+0.093$ ($q=.137$) & $+0.213$ ($q=.002$) \\
1B & referral/disclaimer & $+0.253$ ($q=.002$) & $+0.190$ ($q=.002$) \\
1B & structuring & $+0.337$ ($q=.002$) & $-0.367$ ($q=.002$) \\
\midrule
7B & empathy opener & $+0.263$ ($q=.002$) & $+0.070$ ($q=.387$) \\
7B & validation phrase & $+0.103$ ($q=.121$) & $+0.147$ ($q=.021$) \\
7B & referral/disclaimer & $+0.243$ ($q=.002$) & $+0.083$ ($q=.203$) \\
7B & structuring & $+0.227$ ($q=.002$) & $-0.277$ ($q=.002$) \\
\bottomrule
\end{tabular}%
}
\caption{Change in the advice-minus-factual proxy gap across transitions
(plus-one-corrected prompt-clustered permutation tests, 10k draws; displayed
$q$ is BH-adjusted across all 28 adjacent-stage$\times$proxy tests). The
Base$\to$SFT column also changes the input serialization and is not a clean SFT
effect. Later raw $p\geq.168$.}
\label{tab:interaction}
\end{table}

\paragraph{Comparisons with targeted controls.} The four-seed
Base/SFT/DPO extension tests 24 planned stage$\times$control interactions
(Appendix~\ref{app:matched-stage}). Four survive BH correction at $\tau=.82$.
For empathy relative to emotional non-advice, the 7B Base$\to$SFT/interface
transition increases the advice gap
($+0.358$, $q=.016$), but DPO reverses it at both 1B ($-0.467$, $q=.001$) and
7B ($-0.642$, $q=.001$). At 7B, SFT also shifts structuring toward neutral
advice ($-0.208$, $q=.026$). No validation or referral interaction survives
$q<.05$; 1B DPO validation tendencies have $q=.059$--$.095$. The threshold
sweep yields three significant interactions at $\tau=.80$, four at $.82$, and
five at $.84$; only the two 7B
empathy transitions survive at all three thresholds. Thus the broad
advice-minus-factual changes cannot all be read as advice-specific.

\begin{table*}[t]
\centering
\small
\setlength{\tabcolsep}{4.5pt}
\begin{tabular}{llrrrr}
\toprule
Model & Stage & Empathy & Validation & Referral & Structuring \\
\midrule
1B & base     & .217 (.204) & .073 (.073) & .483 (.430) & .757 (.430) \\
1B & sft      & .393 (.393) & .167 (.167) & .697 (.684) & .847 (.767) \\
1B & dpo      & .550 (.510) & .380 (.380) & .933 (.873) & 1.000 (.400) \\
1B & rlvr     & .567 (.507) & .410 (.410) & .923 (.836) & 1.000 (.440) \\
1B & instruct & .590 (.550) & .363 (.356) & .900 (.820) & 1.000 (.420) \\
\midrule
7B & base     & .230 (.230) & .087 (.080) & .577 (.524) & .853 (.486) \\
7B & sft      & .493 (.493) & .183 (.183) & .767 (.767) & .787 (.714) \\
7B & dpo      & .597 (.564) & .330 (.330) & .923 (.850) & .977 (.437) \\
7B & instruct & .617 (.590) & .383 (.370) & .913 (.853) & .970 (.423) \\
\bottomrule
\end{tabular}
\caption{Stagewise advice detector-positive rate with the \textbf{advice-minus-factual
proxy gap in parentheses} ($\tau=0.82$). Transition tests and corrected values
appear in Table~\ref{tab:interaction}. Base checkpoints use plain QA; all later
checkpoints use chat serialization, so Base$\to$SFT differences conflate stage
and interface.}
\label{tab:stagewise}
\end{table*}

\subsection{DPO phrase differences track response length}
\label{sec:datalayers}

The exact mixtures contain 378{,}301 preference pairs for the 1B model and
378{,}341 for the 7B model. Chosen responses are slightly shorter in aggregate:
the chosen-to-rejected mean-word ratio is .990 for 1B and .997 for 7B. This
aggregate normalization does not remove
within-pair length as a possible covariate. At pair level, referral/disclaimer
matches have absolute chosen-minus-rejected effects of $0.100$ percentage
points (pp; 95\% CI $[.065,.135]$) at 1B and $0.179$ pp
($[.144,.214]$) at 7B; matched odds ratios (MORs) are $1.181$ and $1.352$.
Structuring effects are larger: $1.393$ pp ($[1.278,1.507]$) and $2.280$ pp
($[2.167,2.394]$). Empathy is weakly enriched only at 1B after correction;
validation is not (Table~\ref{tab:dpo}).

Only 19.1\% of 1B pairs and 18.7\% of 7B pairs fall within $\pm20$ words. In
that comparable-length stratum, the referral/disclaimer difference is $-.010$
pp for 1B (MOR .97, raw $p=.763$) and +.009 pp for 7B (MOR 1.03, raw
$p=.810$). The high-precision referral difference is +.017 pp for 1B
($p=.299$) and $-.003$ pp for 7B ($p=.925$). Structuring remains positive but
smaller: +.184 pp for 1B (raw $p=.047$) and +.343 pp for 7B (raw $p=.0002$).
In the noncomparable strata, signs follow the longer response side. Aggregate
DPO enrichment is therefore strongly length-aligned; marginal per-word ratios
do not establish a within-pair preference effect. Conditioning on length also
changes the target quantity and does not identify a causal preference for a
marker.

A narrower referral lexicon that excludes the broad ``consult/talk/speak to a''
phrases remains positive: $0.023$ pp ($[.006,.040]$, MOR $1.17$, raw
$p=.010$) at 1B and $0.043$ pp ($[.026,.061]$, MOR $1.34$,
$p<2\times10^{-6}$) at 7B. The original referral direction is positive in
eight of nine 1B sources and all nine 7B sources; the high-precision direction
is positive in six of nine 1B sources and seven of nine 7B sources. Structuring
is positive in eight of nine sources at each size and is driven
mainly by ``first,'' and ``here are some.'' Only 5.7\% of rows belong to a
duplicated prompt-text group (maximum multiplicity 15), so ordinary row-level
McNemar inference may be mildly optimistic but is not dominated by repeated
prompts. Full phrase/source results and intervals are in
Appendix~\ref{app:dpo-robustness}. These co-occurrences are compatible with,
but do not attribute or directly validate, checkpoint changes.

\begin{table}[t]
\centering
\small
\setlength{\tabcolsep}{2.5pt}
\resizebox{\columnwidth}{!}{%
\begin{tabular}{llrrrrr}
\toprule
Model & Marker proxy & $\Delta$ pp & MOR & Chosen only & Rejected only & BH $q$ \\
\midrule
1B & empathy opener & .035 & 1.082 & 1768 & 1634 & .036 \\
1B & validation phrase & .006 & 1.130 & 183 & 162 & .282 \\
1B & referral/disclaimer & .100 & 1.181 & 2478 & 2098 & $<10^{-7}$ \\
1B & structuring & 1.393 & 1.242 & 27047 & 21779 & $<10^{-124}$ \\
\midrule
7B & empathy opener & .030 & 1.069 & 1734 & 1622 & .074 \\
7B & validation phrase & .009 & 1.224 & 191 & 156 & .078 \\
7B & referral/disclaimer & .179 & 1.352 & 2602 & 1925 & $<10^{-22}$ \\
7B & structuring & 2.280 & 1.435 & 28466 & 19838 & $<10^{-300}$ \\
\bottomrule
\end{tabular}%
}
\caption{Aggregate DPO pair-level effects, unadjusted for within-pair length.
$\Delta$ pp is the absolute paired risk
difference in percentage points; MOR is chosen-only/rejected-only. $q$ is
Benjamini--Hochberg adjusted across eight exact McNemar tests. Confidence
intervals and descriptive per-word ratios are reported in the text and artifact.}
\label{tab:dpo}
\end{table}

\paragraph{Supporting SFT and RLVR searches.}
The assistant-only SFT searches recover seeded phrases, but the comparison is
phrase-sparse and dominated by ``here are some''; per-phrase and leave-one-out
results are in Appendix~\ref{app:perphrase}. The searched RLVR datasets contain
zero or one match for each referral, empathy, and validation phrase, while
generic structuring markers remain. These exact searches establish phrase
co-occurrence in available data, not attribution or a training effect.

\section{Discussion}

\paragraph{What the trajectories measure.}
The checkpoint analysis reveals trajectories of a fixed historical detector,
not semantic advice behaviors. All advice-minus-factual gaps increase across
the confounded Base-to-SFT/interface transition, and same-interface SFT-to-DPO
changes vary across categories. The detector audit shows why the distinction
matters: embedding matches add false positives, and a fixed output window can
reverse a key final-7B contrast. The available evidence therefore supports
comparisons among explicitly named rules and windows, while human-judged
checkpoint trajectories remain open.

\paragraph{Length and preference data.}
Exact matching avoids embedding-threshold errors in the DPO analysis, but it
still gives longer responses more opportunities to contain a phrase. Aggregate
referral/disclaimer and structuring differences largely follow the longer side
of a pair. Among responses within 20 words, referral/disclaimer differences are
near zero and structuring differences are much smaller. Because the length
strata are post hoc and conditioning on length changes the target quantity,
this analysis diagnoses sensitivity rather than a causal preference for or
against any marker.

\paragraph{Implications for measurement.}
Broad historical category names should not be mistaken for narrow constructs.
The referral/disclaimer rule combines recommendations, disclaimers, and
professional-language mentions; the empathy-opener rule searches beyond the
opening sentence. Corpus matches show where the same strings occur, but do not
identify the examples responsible for a model output. Future work should align
category definitions with independent human labels, annotate response prefixes
when evaluating prefix detectors, and cross checkpoint with prompt interface
before making semantic or training-stage claims.

\section{Conclusion}
\label{sec:roadmap}

Conclusions about advice-style language in OLMo-2 change with the detector and
the amount of response text examined. The historical union detector produces
distinct condition and checkpoint trajectories, but its low precision,
interface confounding, and final-7B window sensitivity preclude a semantic
reading. Exact DPO matches are more common on the chosen side in aggregate;
among responses within 20 words, referral/disclaimer differences approach zero
and structuring differences shrink from 1.393--2.280 percentage points to
.184--.343. These results support a focused measurement claim: detection rules
and response length materially alter the observed pattern. They do not show
semantic empathy prevalence, referral quality, or causal effects of training.

\section{Limitations}
\label{sec:limitations}

All analyses are observational. Checkpoint localization, corpus co-occurrence,
and DPO chosen/rejected differences establish neither training-example
attribution nor a causal effect. Base uses plain QA while SFT and later
checkpoints use chat serialization, so every Base-to-SFT contrast combines
stage and interface.

The 120-sentence audit uses AI-consensus labels, has no completed human pass,
and does not validate response-level prevalence. It samples 97 seed-0 responses,
and category-level intervals are wide. Raw-response window checks are available
only for the final 7B checkpoint. They cannot repair or validate trajectories at
the other checkpoints.

The targeted controls are not topic-paired to the 60 advice prompts, and a
non-significant contrast is not evidence of equivalence. Twelve pilot prompts
were reused in the 140-prompt battery; excluding them changes final-7B
advice-minus-factual gaps by at most .045. Five canonical seeds and four
extension seeds do not remove sampling uncertainty.

The DPO length bins are post hoc rather than a pre-specified conditional model,
and 5.7\% of rows belong to duplicated prompt-text groups. SFT recoverability is
based on only three to six distinct three-word types. Stage-2 Dolmino-Mix-1124
was not searched because no public index was available, and the Stage-1 novelty
measure is verbatim-only and length-confounded. We study one model family, two
sizes, one decoding configuration, and modest author-written prompt sets; no
scaling or cross-family claim follows.

\section*{Ethics Statement}

\paragraph{Data and human subjects.} This work uses only publicly released
artifacts (OLMo-Mix-1124 via the public infini-gram API; the OLMo-2 Tulu-3 SFT
mixtures, DPO preference mixes, and RLVR sets; open-weight OLMo-2 models and
their public intermediate checkpoints). We collected no new data from people.
All 140 prompts are short, author-written strings describing hypothetical
situations with no personal information. The study involves no human subjects
and needed no ethics-board review.

\paragraph{Sensitive domain.} Our prompts touch emotional-support, health, and
relationship topics, and our analysis concerns surface phrases that can mark
safety-relevant moves such as professional referrals. The audit sample shows substantial detector
overprediction, suggesting upward-biased absolute rates; we also drop
\texttt{crisis\_referral}. Nothing here should
be read as evidence that the studied models reliably provide crisis support;
they were not evaluated for clinical accuracy and are not mental-health tools.

\paragraph{Dual use.} Making the provenance of a safety-relevant marker inspectable
could in principle guide its removal from a training mixture; we frame this as
a risk to monitor, show co-occurrence rather than causation, and release an
analysis tool, not any model or dataset that could cause direct harm.

\bibliography{custom}

\appendix

\section{Prompts and decoding configuration}
\label{app:prompts}

All answers use nucleus sampling (temperature $0.7$, top-$p$ $0.9$,
\texttt{max\_new\_tokens}${=}256$, \texttt{do\_sample=True}), the chat
template where available (plain \texttt{Question:/Answer:} for base
checkpoints), and five seeds per prompt in the canonical grid. The battery: 60 advice, 30 factual,
20 emotional non-advice, 20 neutral advice, 10 domain factual; the complete
machine-readable list ships as \texttt{experiments/prompts\_full.json}. The
Base/SFT/DPO axis-targeted-control extension uses four seeds; all balanced
extension tests use seeds 0--3 for every condition. Eight advice and four
factual prompts were reused from the earlier 12-prompt pilot; the other 128 are
new to the full battery.

\paragraph{Examples.}
\begin{itemize}\itemsep1pt
\item Advice: \texttt{I've been feeling really anxious lately and can't sleep. Any advice?}
\item Factual: \texttt{How does photosynthesis work?}
\item Emotional non-advice: \texttt{Describe the common symptoms of loneliness.}
\item Neutral advice: \texttt{How should I organize my bookshelf?}
\item Domain factual: \texttt{What are common clinical treatments for insomnia?}
\end{itemize}

\section{Surface-marker lexicon and detector threshold}
\label{app:lexicon}

This section records the executable four-category detector in
\path{experiments/raw_output_robustness.py}, used to reproduce the July
final-7B response analysis. Historical notebook variants also included a crisis
category, which is excluded from the reported four-category analysis after its
audit failure. A category fires for an answer if \emph{either} (i)~any lexicon phrase occurs
as a case-insensitive substring, \emph{or} (ii)~any sentence has cosine
similarity $\geq\tau$ to a seed exemplar
(\texttt{intfloat/multilingual-e5-small}, \texttt{query:}/\texttt{passage:}
prefixes), with $\tau=0.82$ primary. Every sentence is eligible; despite its
historical name, \texttt{empathy\_opener} is not restricted to an opening span.

\paragraph{Lexicon (substring match).}
\begin{itemize}\itemsep1pt
\item \textbf{empathy-opener proxy} (\texttt{empathy\_opener}): \texttt{i'm sorry to hear}, \texttt{i am sorry to hear}, \texttt{i'm really sorry}, \texttt{that sounds really}, \texttt{that sounds very}, \texttt{i can understand}, \texttt{i can imagine}, \texttt{it's understandable}, \texttt{that must be}, \texttt{i'm here for you}, \texttt{i hear you}, \texttt{it sounds like you}
\item \textbf{validation-phrase proxy}: \texttt{your feelings are valid}, \texttt{it's okay to feel}, \texttt{it is okay to feel}, \texttt{it's normal to feel}, \texttt{it's completely normal}, \texttt{what you're feeling}, \texttt{there's nothing wrong with}, \texttt{it's natural to}
\item \textbf{referral/disclaimer proxy} (\texttt{disclaimer}): \texttt{consult a}, \texttt{talk to a}, \texttt{speak with a}, \texttt{see a doctor}, \texttt{seek medical}, \texttt{i am not a doctor}, \texttt{i'm not a doctor}, \texttt{i'm not a medical}, \texttt{medical professional}, \texttt{healthcare professional}, \texttt{mental health professional}, \texttt{seek professional}, \texttt{professional help}
\item \textbf{structuring-marker proxy} (\texttt{structure}): \texttt{here are some}, \texttt{here are a few}, \texttt{consider the following}, \texttt{steps you can take}, \texttt{a few things you can}, \texttt{first,}, \texttt{secondly,}
\end{itemize}

\paragraph{e5 seed exemplars.} Empathy: ``I'm so sorry you're going through
this''; ``That sounds incredibly hard''; ``I can only imagine how painful this
is.'' Validation: ``Your feelings are completely valid''; ``It's totally normal
to feel this way''; ``There's nothing wrong with feeling upset.''
Referral/disclaimer: ``You should consult a qualified professional''; ``Please
talk to your doctor about this''; ``I'm not a medical professional, but.''
Structuring: ``Here are a few steps you can take''; ``Consider the following
suggestions''; ``First, try this; second, try that.''

\section{Threshold sweep}
\label{app:threshold}

Table~\ref{tab:threshold} reports the exact contrasts at $\tau=0.80$, $0.82$,
and $0.84$. Structuring is stable: advice and
neutral-advice remain within $.05$ in every cell. Validation exceeds its
strongest control in all six cells.
Empathy's emotional-minus-advice contrast reverses at $.84$, while referral's
advice-minus-domain contrast changes magnitude. The primary $\tau=.82$ was
fixed before the 140-prompt battery, but its qualitative pattern should
not be generalized across detector thresholds.

\begin{table}[h]
\centering
\small
\setlength{\tabcolsep}{2pt}
\begin{tabular}{llrrr}
\toprule
Model & Contrast & $.80$ & $.82$ & $.84$ \\
\midrule
1B & empathy: emot.$-$adv. & $.030$ & $.200$ & $-.107$ \\
7B &                         & $.030$ & $.143$ & $-.280$ \\
1B & validation: adv.$-$max ctrl. & $.030$ & $.243$ & $.107$ \\
7B &                                & $.047$ & $.253$ & $.170$ \\
1B & referral: adv.$-$domain & $-.003$ & $.080$ & $.093$ \\
7B &                          & $.017$ & $.093$ & $.280$ \\
1B & structuring: adv.$-$neutral & $.000$ & $.000$ & $.047$ \\
7B &                             & $.000$ & $-.030$ & $-.010$ \\
\bottomrule
\end{tabular}
\caption{Threshold robustness of the four decomposition contrasts. Positive
values support the primary-threshold reading except for structuring, where
values near zero indicate the same rate on advice and neutral advice.}
\label{tab:threshold}
\end{table}

\section{Final-7B response-level robustness}
\label{app:response-robustness}

The recovered \texttt{final\_7b\_answers.jsonl} contains all 700 final-7B
responses (140 prompts $\times$ five seeds). Table~\ref{tab:response-length}
shows the opportunity imbalance. ``Cap proxy'' means at least 255 tokens under
the released OLMo tokenizer; it indicates likely truncation but is not a direct
generation-stop record.

\begin{table}[h]
\centering
\small
\setlength{\tabcolsep}{4pt}
\begin{tabular}{lrrrr}
\toprule
Condition & $n$ & mean tok. & mean sent. & cap proxy \\
\midrule
advice & 300 & 243.8 & 13.5 & .880 \\
factual & 150 & 187.4 & 7.8 & .573 \\
emotional & 100 & 253.5 & 10.9 & .980 \\
neutral advice & 100 & 254.6 & 13.8 & .970 \\
domain factual & 50 & 250.5 & 10.2 & .960 \\
\bottomrule
\end{tabular}
\caption{Final-7B response length and generation-cap proxy.}
\label{tab:response-length}
\end{table}

Table~\ref{tab:window-components} recomputes each targeted-control contrast using
exact substring matching alone and the exact-or-embedding union, over full
answers and fixed prefix windows. The full-union column exactly reproduces the
fixed final-7B result. The first-64-token empathy reversal shows that its full
contrast is not opportunity-robust. Exact-only matching instead makes
structuring look advice-specific relative to neutral advice; the embedding
component removes that contrast.

\begin{table}[h]
\centering
\small
\setlength{\tabcolsep}{3pt}
\resizebox{\columnwidth}{!}{%
\begin{tabular}{lrrrr}
\toprule
Proxy contrast & Exact full & Union full & Union 64 & Union 128 \\
\midrule
empathy: adv.$-$emot. & .323 & $-.147$ & .253 & .063 \\
validation: adv.$-$emot. & .153 & .250 & .177 & .210 \\
referral: adv.$-$domain & .157 & .090 & .113 & .117 \\
structuring: adv.$-$neutral & .390 & $-.030$ & $-.033$ & $-.030$ \\
\bottomrule
\end{tabular}%
}
\caption{Detector-component and fixed-window sensitivity at final 7B
($\tau=.82$ for union columns).}
\label{tab:window-components}
\end{table}

Among advice responses, union detections include one category in 4.7\% of
cases, two in 29.0\%, three in 40.7\%, and four in 25.7\%. Pairwise phi
correlations range from $-.14$ to $.21$; referral/disclaimer appears in 93.0\%
of validation-positive responses. Excluding the eight advice and four factual
pilot prompts changes the empathy advice-minus-factual gap from .587 to .542,
validation from .367 to .369, referral/disclaimer from .850 to .831, and
structuring from .423 to .458.

\section{Checkpoint interactions with targeted controls}
\label{app:matched-stage}

Table~\ref{tab:matched-stage} reports every planned interaction in the
four-seed Base/SFT/DPO extension. Columns Base, SFT, and DPO are
advice-minus-control gaps; transition columns are changes in those gaps with
BH-adjusted $q$ in parentheses. The family contains all 24 displayed
transition tests.

\begin{table*}[h]
\centering
\footnotesize
\setlength{\tabcolsep}{2.5pt}
\begin{tabular}{lllrrrrr}
\toprule
Model & Marker & Control & Base & SFT & DPO & Base$\to$SFT & SFT$\to$DPO \\
\midrule
1B & referral & domain factual & .188 & .308 & .167 & $+.121$ (.414) & $-.142$ (.296) \\
1B & empathy & emotional & .017 & .225 & $-.242$ & $+.208$ (.095) & $-.467$ (.001) \\
1B & structuring & neutral advice & $-.013$ & $-.050$ & .000 & $-.037$ (.653) & $+.050$ (.475) \\
1B & validation & domain factual & .071 & .146 & .375 & $+.075$ (.475) & $+.229$ (.095) \\
1B & validation & emotional & .021 & .121 & .225 & $+.100$ (.293) & $+.104$ (.296) \\
1B & validation & neutral advice & .058 & .171 & .350 & $+.113$ (.189) & $+.179$ (.059) \\
\midrule
7B & referral & domain factual & .321 & .067 & .200 & $-.254$ (.166) & $+.133$ (.293) \\
7B & empathy & emotional & .108 & .467 & $-.175$ & $+.358$ (.016) & $-.642$ (.001) \\
7B & structuring & neutral advice & .050 & $-.158$ & $-.025$ & $-.208$ (.026) & $+.133$ (.095) \\
7B & validation & domain factual & .092 & .175 & .317 & $+.083$ (.475) & $+.142$ (.296) \\
7B & validation & emotional & .054 & .125 & .179 & $+.071$ (.414) & $+.054$ (.537) \\
7B & validation & neutral advice & .092 & .163 & .292 & $+.071$ (.414) & $+.129$ (.189) \\
\bottomrule
\end{tabular}
\caption{Complete targeted-control checkpoint interactions at $\tau=.82$ (10k
prompt-clustered permutations). Four of 24 transition tests have $q<.05$.
Only the 7B empathy Base$\to$SFT and SFT$\to$DPO effects retain $q<.05$ at
all three detector thresholds.}
\label{tab:matched-stage}
\end{table*}

\section{Per-phrase recoverability and leave-one-out}
\label{app:perphrase}

Per-phrase assistant-only counts against each model's exact mixture
(\texttt{-0225}, $866{,}138$ docs, 1B; \texttt{tulu-3-sft-olmo-2-mixture},
$939{,}344$, 7B); \texttt{infl} $=$ blob/assistant.

\begin{table}[h]
\centering
\small
\setlength{\tabcolsep}{4pt}
\resizebox{0.88\columnwidth}{!}{%
\begin{tabular}{llrrr}
\toprule
Model & Phrase & Asst. & Blob & Infl. \\
\midrule
1B & first,                     & 44422 & 52676 & 1.19 \\
1B & here are some              & 11914 & 13276 & 1.11 \\
1B & here are a few             & 5421  & 6056  & 1.12 \\
1B & healthcare professional    & 3075  & 3216  & 1.05 \\
1B & mental health professional & 2109  & 2186  & 1.04 \\
1B & professional help          & 1765  & 1825  & 1.03 \\
1B & seek professional          & 1353  & 1377  & 1.02 \\
1B & it sounds like you         & 1130  & 1164  & 1.03 \\
1B & consult a                  & 891   & 978   & 1.10 \\
1B & it's understandable       & 218   & 227   & 1.04 \\
1B & it's natural to            & 152   & 157   & 1.03 \\
\midrule
7B & here are some              & 12049 & 13432 & 1.11 \\
7B & here are a few             & 5511  & 6156  & 1.12 \\
7B & healthcare professional    & 3106  & 3309  & 1.07 \\
7B & i'm really sorry           & 2369  & 2378  & 1.00 \\
7B & mental health professional & 2120  & 2214  & 1.04 \\
7B & professional help          & 1774  & 1838  & 1.04 \\
7B & seek professional          & 1355  & 1379  & 1.02 \\
7B & it sounds like you         & 1139  & 1184  & 1.04 \\
7B & consult a                  & 927   & 1031  & 1.11 \\
7B & talk to a                  & 286   & 408   & 1.43 \\
7B & it's completely normal     & 78    & 103   & 1.32 \\
\bottomrule
\end{tabular}%
}
\caption{Per-phrase SFT-assistant recoverability (representative rows). The 3
(1B) / 6 (7B) distinct 3-word types drive Section~\ref{sec:datalayers};
``here are some'' dominates the 3-word mean.}
\label{tab:perphrase}
\end{table}

\paragraph{Length-matched buckets (occurrence-weighted).} Seeded-marker vs
topical per-million for 1B: 2-word $1884.3$ vs $1760.2$ ($1.1\times$, $n_{\text{occ}}{=}14$),
3-word $12900.2$ vs $629.8$ ($20.5\times$, $45$), 4-word $1596.9$ vs $11.7$
($136.9\times$, $30$). For 7B, 2-word $1784.5$ vs $996.8$ ($1.8\times$, $20$),
3-word $9354.4$ vs $1280.4$ ($7.3\times$, $55$), 4-word $1396.6$ vs $96.2$
($14.5\times$, $24$). \textbf{Distinct-phrase leave-one-out (3-word):}
dropping ``here are some'' moves the mean from $5455.2$ to $1305.2$/M at 1B
(3 types) and from $3026.0$ to $1065.9$/M at 7B (6 types).

\section{DPO phrase and source robustness}
\label{app:dpo-robustness}

Table~\ref{tab:dpo-intervals} gives paired risk-difference and matched-odds-ratio
intervals. The higher-precision referral set requires an explicit doctor or
professional object (e.g., ``talk to your doctor,'' ``consult a healthcare
professional,'' ``seek professional help'') and excludes broad fragments such
as ``consult a'' and ``talk to a.'' Intervals use the discordant-pair normal
approximation for the risk difference and log-MOR; exact McNemar $p$-values are
released.

\begin{table*}[h]
\centering
\small
\setlength{\tabcolsep}{4pt}
\begin{tabular}{lllrrr}
\toprule
Model & Phrase group & $\Delta$ pp [95\% CI] & MOR [95\% CI] & C-only & R-only \\
\midrule
1B & referral/disclaimer & .100 [.065,.135] & 1.181 [1.114,1.252] & 2478 & 2098 \\
1B & high-precision referral & .023 [.006,.040] & 1.170 [1.039,1.318] & 591 & 505 \\
1B & structuring & 1.393 [1.278,1.507] & 1.242 [1.220,1.264] & 27047 & 21779 \\
\midrule
7B & referral/disclaimer & .179 [.144,.214] & 1.352 [1.274,1.434] & 2602 & 1925 \\
7B & high-precision referral & .043 [.026,.061] & 1.335 [1.187,1.502] & 649 & 486 \\
7B & structuring & 2.280 [2.167,2.394] & 1.435 [1.409,1.461] & 28466 & 19838 \\
\bottomrule
\end{tabular}
\caption{DPO paired absolute and relative effects. Positive $\Delta$ means more
chosen-only than rejected-only matches.}
\label{tab:dpo-intervals}
\end{table*}

Table~\ref{tab:dpo-length} is the requested within-pair length diagnostic. The
near-antisymmetric signs in the longer-response strata show that the aggregate
chosen-side effect is strongly associated with relative response length. The
comparable-length band is descriptive and post hoc; it is not a substitute for
a pre-specified conditional model.

\begin{table*}[h]
\centering
\footnotesize
\setlength{\tabcolsep}{3pt}
\begin{tabular}{llrrrrrrr}
\toprule
Model & Word-difference band & $n$ & \multicolumn{2}{c}{Referral} &
\multicolumn{2}{c}{High precision} & \multicolumn{2}{c}{Structuring} \\
 & & & $\Delta$ pp & MOR & $\Delta$ pp & MOR & $\Delta$ pp & MOR \\
\midrule
1B & chosen longer $>20$ & 171170 & +.470 & 1.991 & +.059 & 1.408 & +5.979 & 2.361 \\
1B & rejected longer $>20$ & 134707 & $-.310$ & .616 & $-.020$ & .871 & $-3.786$ & .579 \\
1B & within $\pm20$ & 72424 & $-.010$ & .965 & +.017 & 1.238 & +.184 & 1.062 \\
\midrule
7B & chosen longer $>20$ & 187201 & +.502 & 2.172 & +.098 & 1.761 & +6.562 & 2.610 \\
7B & rejected longer $>20$ & 120381 & $-.223$ & .706 & $-.015$ & .905 & $-3.239$ & .619 \\
7B & within $\pm20$ & 70759 & +.009 & 1.028 & $-.003$ & .966 & +.343 & 1.123 \\
\bottomrule
\end{tabular}
\caption{DPO effects stratified by chosen-minus-rejected response length.
Positive $\Delta$ means more chosen-only matches.}
\label{tab:dpo-length}
\end{table*}

For the original referral group, positive chosen-minus-rejected differences
occur in eight of nine 1B sources (one tie) and all nine 7B sources. Under the
higher-precision group they occur in six of nine 1B sources and seven of nine
7B sources. Structuring is positive in eight of nine sources at each size; the
negative source is the instruction-following personas subset. The main
structuring contribution comes from ``first,'' (.931 pp for 1B and 1.381 pp for
7B) and ``here are some'' (.339 pp for 1B and .736 pp for 7B), rather than every
phrase moving uniformly. The main referral contributors are ``professional
help,'' ``consult a,'' and named professional phrases. There are 367{,}082
distinct 1B prompt texts and 367{,}110 distinct 7B prompt texts; 21{,}551 1B
rows and 21{,}574 7B rows (5.7\% for each) belong to duplicated groups, with maximum
multiplicity 15.

\section{Response-cluster audit uncertainty}
\label{app:audit-bootstrap}

The 120 audit sentences come from 97 seed-0 responses. To preserve dependence
among sentences from the same response, we assign Dirichlet$(1)$ weights to
response keys in a 10{,}000-draw Bayesian bootstrap and recompute each metric.
Table~\ref{tab:audit-bootstrap} reports percentile intervals. These intervals
quantify sampling uncertainty around the AI consensus labels; they do not turn
the audit into human validation.

\begin{table}[t]
\centering
\footnotesize
\setlength{\tabcolsep}{2pt}
\begin{tabular}{lccc}
\toprule
Category & Precision 95\% & Recall 95\% & F1 95\% \\
\midrule
empathy expression & [.04,.46] & [.10,.90] & [.05,.57] \\
validation         & [.02,.84] & [.03,.98] & [.02,.83] \\
referral           & [.00,.22] & [1.00,1.00] & [.00,.36] \\
structuring        & [.35,.63] & [.74,.96] & [.49,.74] \\
\bottomrule
\end{tabular}
\caption{Bayesian response-cluster bootstrap intervals for the 120-sentence
AI consensus audit. Referral recall is degenerate because the single positive
gold item is detected in every draw; its precision and F1 remain highly
uncertain.}
\label{tab:audit-bootstrap}
\end{table}

\section{Statistical formulas}
\label{app:formulas}

Primary inference uses two-sided prompt-clustered permutation tests on
prompt-level means (10{,}000 draws), with Monte Carlo
$p=(b+1)/(B+1)$; $p<10^{-4}$ denotes zero raw exceedances.
The distinct-phrase leave-one-out removes the most frequent 3-word phrase before
recomputing the mean seeded-marker per-million. DPO comparisons use exact McNemar
tests on chosen-only versus rejected-only pairs, Haldane--Anscombe matched odds
ratios with log-scale intervals, paired risk differences with normal-approximation
intervals, and phrase occurrences per million whitespace-delimited response
words. Released scripts derive all statistics from the artifact.

\end{document}
